\begin{table}[htbp] 
	\centering
\caption{Covid-19 incidence rates and state prime ministers' statements favoring (de)centralized decision-making}
 \label{tb.mainreg}
\begin{tabular}{lccccc} \toprule
    & (1) & (2) & (3) & (4) & (5) \\ \midrule
ln(state-level  & 1.730** & 1.904*** & 2.003*** & 1.983** & 1.985** \\
\, incidence rate)	& [1.22,2.45] & [1.34,2.70]  & [1.41,2.85]  & [1.30,3.02]  & [1.29,3.06] \\
 & & & & & \\
ln(state-level  &  &  &  &  & 0.806 \\
\, vacc. rate)		&  &  &  &  & [0.606,1.072] \\
		& & & & & \\
Linear time trend & No & Yes & Yes & Yes & Yes \\
		& & & & & \\
Dummies Eastern  &    &    &     &     & \\
\,  states and Berlin  & No & No & Yes & Yes & Yes \\ 
		& & & & & \\
State-level controls & No & No & No & Yes & Yes \\ \\ \midrule
 Observations & 600 & 600 & 600 & 600 & 600 \\		\bottomrule
\end{tabular}
\label{tb.reg01}
\begin{minipage}{\textwidth}
	\small
\textit{Notes:} Each column reports exponentiated coefficients from a separate ordered logit regression. Dependent variable is the 4-point pandemic policy index (PPI). Higher values indicate a stronger expressed desire for centralized decision-making. Coefficients are expressed as changes in the odds ratio of giving a statement $Y$ that falls into a higher vs a lower category $j$, i.e., $P(Y\leq j)/P(Y > j)$. All regressions control for ln(federal level incidence rate) and column (5) additionally controls for ln(federal level vaccination rate) and a dummy for the pre-vaccination period. State-level controls include states' general stances towards federalism, share of the manufacturing sector in gross value added, and financial strength ("Finanzkraftmesszahl"). \newline \textit{*** $p < 0.01$, ** $p < 0.05$, * $p < 0.10$}
	\end{minipage}
\end{table}